\documentclass[11pt]{article}

\usepackage[margin=1in]{geometry}
\usepackage{amsmath}
\usepackage{amssymb}
\usepackage{booktabs}
\usepackage{graphicx}
\usepackage{hyperref}
\usepackage{caption}
\usepackage{array}
\usepackage{xcolor}
\usepackage{enumitem}

\title{\textbf{A Compact CNN-LSTM Network Intrusion Detection System \\
with Leave-One-Attack-Out Generalization Analysis on CICIDS2017}}
\author{Gopal Deswal \\ \small Under the guidance of: Sandeep Gupta}
\date{September 2026}

\begin{document}
\maketitle

\begin{abstract}
Network Intrusion Detection Systems (NIDS) based on deep learning have reported near-ceiling
accuracies on the CICIDS2017 benchmark, with published CNN-LSTM and hybrid architectures
routinely exceeding 99\% under conventional random-split evaluation. We argue, and empirically
demonstrate, that this evaluation protocol is a poor proxy for real-world robustness: a model can
achieve arbitrarily high accuracy while having learned to recognize only the specific attack
categories present at training time, with no guarantee of behavior on a category it has never
encountered -- precisely the scenario a zero-day attack presents. In this work we first build a
compact CNN-LSTM classifier (33{,}761 parameters, roughly an order of magnitude smaller than
typical published architectures) that reaches 98.29\% test accuracy and 0.9989 ROC-AUC under
standard evaluation, with a train-test accuracy gap of 0.0005 and a 5-fold cross-validation
standard deviation of 1.39\%, indicating low bias and low variance under this protocol. We then
subject the same model to a stricter Leave-One-Attack-Out (LOAO) evaluation, in which an entire
attack category is withheld from training and treated as unseen at test time, across four
categories (DoS Hulk, PortScan, DDoS, Bot). Accuracy falls to 63--80\% depending on category, with
a complete recall failure (0.00\%) on Bot traffic -- a failure mode that, notably, exactly matches
an independent finding by Xu and Liu (2025) on a different CNN architecture and category set,
suggesting a structural rather than implementation-specific cause. We provide a distribution-shift
argument, grounded in empirical risk minimization theory, for why this collapse is expected rather
than anomalous. We then test whether two classifier-independent signals -- embedding-space distance
to class centroids and benign-only autoencoder reconstruction error -- can recover any of this lost
performance, via a formal four-condition ablation. We find negligible improvement (under 0.3
percentage points in every configuration), a negative result we use to localize the true bottleneck
to the learned representation itself rather than to the decision rule applied on top of it. On
standard evaluation our model is competitive with, though does not exceed, the strongest published
figures (99.49--99.97\%), while using substantially fewer parameters; under LOAO evaluation, to our
knowledge no compared work reports this protocol at all, making direct superiority claims
inapplicable rather than favorable or unfavorable. We conclude with a concrete, mathematically
motivated direction for closing this gap: replacing centroid comparison with instance-level
retrieval and training the embedding space itself under a metric-learning objective rather than
plain cross-entropy classification.
\end{abstract}

\tableofcontents
\newpage

\section{Introduction}

\subsection{Motivation}

Network intrusion detection sits at the front line of modern cybersecurity infrastructure, tasked
with distinguishing malicious traffic from the overwhelming volume of benign network activity in
real time. The rise of deep learning has reshaped this field: hybrid architectures combining
convolutional layers (for local statistical feature extraction across flow-level attributes) with
recurrent layers such as LSTM (for compositional feature aggregation across a flow's temporal
structure) have become a dominant paradigm for NIDS research on the CICIDS2017 benchmark, with
published results frequently exceeding 99\% accuracy \cite{afraji2025,kumar2023}. On the surface,
this suggests the problem of network intrusion detection is close to solved.

We argue this conclusion is premature, and that the near-ceiling accuracy figures dominating the
literature reflect a property of the evaluation protocol rather than of the underlying detection
capability. Under a standard random train/test split, every attack category present in the test
set was also present, in different specific instances, during training. A model achieving 99\%
accuracy under this protocol has demonstrated that it can generalize \emph{within} the distribution
of attack behaviors it was shown -- a meaningfully narrower claim than generalizing to \emph{novel}
attack behavior, which is the actual threat model a deployed NIDS must contend with. A zero-day
attack is, by definition, a category absent from any training corpus available at the time the
model was built. No amount of accuracy on known categories logically implies anything about
behavior on an unknown one; this is the central methodological gap this work is built to probe.

\subsection{Problem Formalization}
\label{sec:formalization}

Let $f_\theta : \mathcal{X} \to [0,1]$ denote the trained classifier, parameterized by $\theta$,
fit via empirical risk minimization over a training set $\{(x_i, y_i)\}_{i=1}^{n}$ drawn i.i.d.\
from a joint distribution $P_{\text{train}}(x, y)$:

\begin{equation}
\theta^{*} = \arg\min_{\theta} \; \frac{1}{n} \sum_{i=1}^{n} \mathcal{L}\big(f_\theta(x_i), y_i\big)
\;+\; \lambda_1 \lVert \theta \rVert_1 \;+\; \lambda_2 \lVert \theta \rVert_2^2
\label{eq:erm}
\end{equation}

where $\mathcal{L}$ is binary cross-entropy and the $\ell_1/\ell_2$ penalty terms constrain the
effective hypothesis space, directly trading a small increase in bias for a reduction in variance
-- the mechanism underlying the low train--test gap reported in Section~\ref{sec:standard-results}.
Standard generalization guarantees for this procedure (e.g., PAC-style bounds) hold under the
assumption that test examples are drawn from the same distribution $P_{\text{train}}$ as the
training data. Under conventional random-split evaluation, this assumption is satisfied by
construction: every category contributing to the test distribution also contributed, independently,
to the training distribution.

The Leave-One-Attack-Out (LOAO) protocol we adopt deliberately violates this assumption. Let
$\mathcal{C} = \{c_1, \ldots, c_k\}$ denote the set of attack categories, and let
$P_{c_j}(x \mid y=1)$ denote the conditional distribution of flow features for category $c_j$.
Under LOAO, the model is trained on $P_{\text{train}}(x, y \mid \text{category} \neq c_j)$ for some
withheld $c_j$, and evaluated on test examples drawn from $P_{c_j}(x \mid y=1)$ -- a distribution
whose support the model has never observed under the positive class during training. No portion of
the standard ERM generalization argument applies to this regime; the model's behavior on
$P_{c_j}$ is, formally, extrapolation outside its training support, with no theoretical guarantee
of correctness. We treat the resulting accuracy collapse (Section~\ref{sec:loao-results}) not as an
anomaly requiring explanation away, but as the \emph{expected} consequence of this distribution
mismatch -- and we use this framing to motivate why our subsequent augmentation attempt
(Section~\ref{sec:augmentation}) targets the representation itself, since distance-based or
reconstruction-based scoring rules computed within an embedding space shaped entirely by
$P_{\text{train}}$ inherit the same support limitation. An embedding function
$g_\theta : \mathcal{X} \to \mathbb{R}^{d}$ (here, the model's penultimate 32-dimensional layer) is
itself a learned function of $\theta^{*}$, and $\theta^{*}$ was optimized with no exposure to
$P_{c_j}$; there is no reason, a priori, to expect distances computed in $g_\theta(\mathcal{X})$ to
separate $P_{c_j}$ from $P_{\text{benign}}$ any more reliably than $f_\theta$ itself does.

\subsection{Contributions}

This work makes the following contributions:

\begin{enumerate}
\item A compact, parameter-efficient CNN-LSTM classifier (33{,}761 parameters) evaluated under the
standard CICIDS2017 protocol, with full bias-variance diagnostics (5-fold cross-validation,
train/validation/test gap analysis) and a comparison against published state-of-the-art results.
\item A Leave-One-Attack-Out evaluation across four withheld categories, quantifying the
generalization gap between standard and zero-day-style evaluation that is absent from the compared
literature \cite{afraji2025,kumar2023}.
\item A formal, component-wise ablation testing whether embedding-distance and reconstruction-error
signals can recover LOAO performance, reported as a negative result with an explicit
distribution-shift explanation for why this outcome is theoretically expected rather than a failure
of implementation.
\item A cross-study convergence check against an independently reported CNN zero-recall failure on
Bot traffic \cite{xu2025}, strengthening the claim that this specific failure mode is structural to
the dataset and supervised paradigm rather than an artifact of either individual implementation.
\end{enumerate}

\subsection{Summary of Findings Relative to Prior Work}

To state plainly where this work stands relative to the literature it compares against: under
standard evaluation, our model does not exceed the highest published accuracy (99.97\%,
\cite{kumar2023}) or the strongest published F1 (0.9956, \cite{afraji2025}); it performs comparably,
within the same 97--99.9\% tier reported across CICIDS2017 studies, while using roughly an order of
magnitude fewer parameters than either compared architecture, neither of which reports its
parameter count. Under LOAO evaluation, direct comparison is not possible in the strict sense, since
neither \cite{afraji2025} nor \cite{kumar2023} report results under this protocol at all -- our
contribution here is introducing this harder evaluation to a setting where it has not previously
been applied, rather than outperforming a directly comparable number. Where a LOAO-style comparison
\emph{is} available, against Xu and Liu's unknown-attack evaluation \cite{xu2025}, our aggregate
LOAO accuracy (70.66\%) is broadly comparable in magnitude to their reported figure (58.82\%), and
our independent replication of their exact 0.0000 Bot-recall failure is, we argue, a more valuable
finding than a favorable accuracy comparison would be, since it is evidence of a shared, structural
limitation rather than an implementation-specific weakness in either study. We emphasize that we do
not claim to outperform any individually comparable published figure, category-for-category; where
a direct category-level comparison exists (Section~\ref{sec:related}), our result is either behind
it or matches a known failure mode. Our augmentation ablation, finally, is an explicit negative
result: neither centroid-distance nor reconstruction-error signals improve on the raw classifier's
baseline LOAO performance, which we argue is itself an informative finding about where the true
bottleneck in zero-day NIDS generalization lies.

\subsection{Paper Organization}

Section~\ref{sec:related} reviews related work on CICIDS2017 classification and
zero-day/unseen-attack detection. Section~\ref{sec:dataset} describes the dataset and preprocessing
pipeline. Section~\ref{sec:architecture} details the proposed architecture. Section~\ref{sec:standard-results}
reports standard-evaluation results and the comparative analysis against published state-of-the-art
methods. Section~\ref{sec:loao} presents the core contribution of this work: the LOAO evaluation,
the augmentation ablation, and a mathematically grounded discussion of the negative result.
Section~\ref{sec:deployment} discusses deployment implications, and Section~\ref{sec:conclusion}
concludes with limitations and concrete directions for future work.

\section{Related Work}
\label{sec:related}

\subsection{Supervised CNN-LSTM Architectures on CICIDS2017}

Several published CNN-LSTM and hybrid architectures report near-ceiling accuracy on CICIDS2017
under standard evaluation. Afraji et al.\ \cite{afraji2025} propose a parallel-fusion architecture
that extracts spatial features via CNN, short-term temporal features via GRU, and long-term
temporal dependencies via LSTM in parallel branches before fusion, achieving 99.49\% accuracy,
99.25\% precision, and 99.88\% recall on CICIDS2017 -- the strongest published F1-score
(0.9956) among the works surveyed for this comparison. Kumar and Sharma \cite{kumar2023} instead
emphasize feature engineering: an intelligent hybrid feature selection stage reduces the input
space to 27 discriminative features before classification, reaching 99.967\% accuracy, the highest
published figure surveyed here. Neither work reports performance under a held-out-category
(zero-day) evaluation protocol; both are evaluated exclusively under random-split conditions where
every attack category represented in the test set is also represented, by other instances, in the
training set.

\subsection{Autoencoder-Based and Unsupervised Zero-Day Detection}

A separate line of work sidesteps the held-out-category problem by never training on attack labels
at all. Hindy et al.\ \cite{hindy2020} train an autoencoder exclusively on benign traffic and use
reconstruction error, thresholded at several operating points (0.05, 0.10, 0.15), as an anomaly
signal to flag zero-day attacks. Their per-category detection accuracy on CICIDS2017 varies
substantially by attack type: attacks with large statistical separation from benign traffic (e.g.,
DoS Hulk, DDoS) are detected with 92--99\% accuracy even at conservative thresholds, whereas attacks
with subtler, benign-mimicking behavior (e.g., Botnet, PortScan at strict thresholds) are detected
far less reliably, in some configurations below 20\% before threshold relaxation. This
category-dependent difficulty is a recurring theme that our own LOAO results reproduce
independently under a supervised rather than unsupervised paradigm, suggesting the difficulty is a
property of the attack categories themselves rather than of any single detection methodology.

This general finding -- that anomaly-detection paradigms trained without attack labels can
sometimes generalize better to unseen threats than purely supervised classifiers, at the cost of
elevated false-alarm rates -- is consistent with the broader anomaly detection literature
\cite{chandola2009}, and with evaluation methodologies developed for related network security
benchmarks such as UNSW-NB15 \cite{moustafa2016}.

\subsection{Controlled Known-vs-Unknown Attack Comparisons}

Most directly relevant to the present work, Xu and Liu \cite{xu2025} conduct a controlled comparison
of supervised (MLP, CNN) and unsupervised (OCSVM, LOF) models on CICIDS2017 under both a standard
evaluation and an explicit unknown-attack test set (DoS slowloris, DoS Slowhttptest, and Bot
withheld from training). Their supervised CNN model achieves 96.50\% accuracy on the standard test
set but collapses to 58.82\% accuracy and 19.54\% recall on the unknown-attack set. Critically, their
reported per-class accuracy for CNN on the withheld Bot category is exactly $0.0000$ -- an identical
total failure mode to what we independently observe for Bot in our own LOAO evaluation
(Section~\ref{sec:loao-results}). This convergent finding across two independent studies, on the
same dataset, using different CNN architectures and different withheld category sets, is strong
evidence that the Bot-detection failure is a structural property of the dataset and supervised
learning paradigm, not an artifact of a single implementation. We discuss the theoretical basis for
why such convergence is expected, rather than coincidental, in Section~\ref{sec:discussion}.

\subsection{Positioning of This Work}

The literature reviewed above splits into two largely non-overlapping camps: works that push
standard-evaluation accuracy toward its ceiling without testing generalization to unseen categories
\cite{afraji2025,kumar2023}, and works that study unseen-category generalization but do not combine
it with a rigorous component-wise ablation of candidate remedies
\cite{hindy2020,xu2025}. This work sits at the intersection: it builds a standard-evaluation-competitive
classifier, subjects it to the harder protocol, and -- rather than stopping at reporting the
collapse -- formally tests, with an isolated ablation rather than a single before/after number,
whether two plausible remedies recover any of the lost performance.

\section{Dataset and Preprocessing}
\label{sec:dataset}

We use the CICIDS2017 dataset, comprising 2{,}830{,}743 labeled network flow records across 15
classes (1 benign, 14 attack sub-types), generated by the Canadian Institute for Cybersecurity to
reflect realistic enterprise network traffic. The raw class distribution is severely imbalanced:
Benign traffic accounts for 85.46\% of records (1{,}977{,}318 samples), while the rarest class,
Heartbleed, contains only 11 samples -- an imbalance ratio of approximately 179{,}756:1. The full
label distribution is summarized in Table~\ref{tab:classdist}.

\begin{table}[h]
\centering
\caption{CICIDS2017 class distribution (top classes by frequency).}
\label{tab:classdist}
\begin{tabular}{lrr}
\toprule
Class & Count & \% of Total \\
\midrule
Benign & 1{,}977{,}318 & 85.457 \\
DoS Hulk & 172{,}846 & 7.470 \\
DDoS & 128{,}014 & 5.533 \\
DoS GoldenEye & 10{,}286 & 0.445 \\
FTP-Patator & 5{,}931 & 0.256 \\
DoS slowloris & 5{,}385 & 0.233 \\
DoS Slowhttptest & 5{,}228 & 0.226 \\
SSH-Patator & 3{,}219 & 0.139 \\
PortScan & 1{,}956 & 0.085 \\
Bot & 1{,}437 & 0.062 \\
Heartbleed & 11 & 0.000 \\
\bottomrule
\end{tabular}
\end{table}

\subsection{Preprocessing Pipeline}

For standard classification, the 14 attack sub-types are collapsed to a single ``Attack'' label,
yielding a binary Benign/Attack task, and the dataset is balanced-sampled for computational
tractability. All flow-level features are standardized to zero mean and unit variance using
\texttt{StandardScaler}, fit exclusively on the training partition to prevent information leakage
from validation or test distributions into the scaling statistics:

\begin{equation}
\hat{x}_{i,k} = \frac{x_{i,k} - \mu_k^{\text{train}}}{\sigma_k^{\text{train}}}
\end{equation}

where $\mu_k^{\text{train}}$ and $\sigma_k^{\text{train}}$ are the mean and standard deviation of
feature $k$ computed only over training samples. Standardized features are reshaped to a
(features, 1) sequence format to serve as input to the Conv1D front end. A stratified 70/15/15
train/validation/test split is used for the standard-evaluation experiments, and an independent
5-fold stratified cross-validation is run separately for the bias-variance diagnostic reported in
Section~\ref{sec:standard-results}.

\subsection{LOAO-Specific Partitioning}

For the Leave-One-Attack-Out experiments (Section~\ref{sec:loao}), a separate partitioning
procedure is required. For a given withheld category $c_j$, all records labeled $c_j$ are removed
from the training and validation pools entirely; the training set retains benign traffic and every
\emph{other} attack category. The test set is constructed from held-out benign traffic plus every
record of the withheld category $c_j$, so that the test distribution for the positive class is
$P_{c_j}(x \mid y=1)$ exactly, with no leakage of $c_j$ instances into any portion of training. This
partitioning is repeated independently for each of the four withheld categories evaluated in this
work (DoS Hulk, PortScan, DDoS, Bot), with a fresh model trained from random initialization for each.

\section{Proposed Architecture}
\label{sec:architecture}

The classifier is a compact hybrid CNN-LSTM with two convolutional blocks feeding a single LSTM
layer, summarized in Table~\ref{tab:arch}.

\begin{table}[h]
\centering
\caption{Proposed CNN-LSTM architecture.}
\label{tab:arch}
\begin{tabular}{l}
\toprule
\textbf{Layer} \\
\midrule
Conv1D(64 filters, kernel=3, ReLU, L1L2 regularized) \\
BatchNormalization \\
MaxPooling1D(pool=2) \\
Conv1D(32 filters, kernel=3, ReLU, L1L2 regularized) \\
BatchNormalization \\
MaxPooling1D(pool=2) \\
LSTM(64 units) \\
Dropout(0.3) \\
Dense(32, ReLU) \\
Dropout(0.3) \\
Dense(1, sigmoid) \\
\bottomrule
\end{tabular}
\end{table}

\subsection{Layer-Level Formulation}

Each one-dimensional convolution computes, for input sequence $x$, kernel $w$ of size $K$, and bias
$b$:

\begin{equation}
(x * w)[t] = \sum_{k=0}^{K-1} x[t+k]\, w[k] + b, \qquad \text{followed by } \text{ReLU}(z) = \max(0, z).
\end{equation}

Each max-pooling operation reduces resolution and introduces local translation invariance:

\begin{equation}
y[t] = \max_{0 \le k < P} z[t \cdot S + k],
\end{equation}

where $P$ is the pooling window and $S$ the stride. The LSTM layer updates a hidden state
$h_t$ and cell state $c_t$ via the standard gated recurrence:

\begin{align}
f_t &= \sigma(W_f [h_{t-1}, x_t] + b_f) \\
i_t &= \sigma(W_i [h_{t-1}, x_t] + b_i) \\
\tilde{c}_t &= \tanh(W_c [h_{t-1}, x_t] + b_c) \\
c_t &= f_t \odot c_{t-1} + i_t \odot \tilde{c}_t \\
o_t &= \sigma(W_o [h_{t-1}, x_t] + b_o) \\
h_t &= o_t \odot \tanh(c_t)
\end{align}

which gives the recurrent layer the capacity to aggregate information across the convolved feature
sequence while selectively retaining or discarding state, a property the CNN front end alone lacks.
The final hidden state is passed through Dropout, a Dense(32) bottleneck -- which also serves as the
embedding used for the distance-based augmentation in Section~\ref{sec:augmentation} -- a second
Dropout, and a final sigmoid unit for binary classification:

\begin{equation}
\hat{y} = \sigma\big(W_{\text{out}} \cdot \text{Dense}_{32}(h_T) + b_{\text{out}}\big) = \frac{1}{1 + e^{-z}}.
\end{equation}

\subsection{Regularization and Parameter Budget}

L1/L2 regularization ($\lambda_1 = 10^{-5}$, $\lambda_2 = 10^{-4}$ in Equation~\ref{eq:erm}) is
applied to convolutional kernels to control overfitting without materially constraining model
capacity. The model totals \textbf{33{,}761 parameters} (33{,}569 trainable), an order of magnitude
smaller than most published CICIDS2017 architectures, which typically do not report parameter
counts at all -- a property we revisit in Section~\ref{sec:deployment} as a deployment advantage
independent of raw accuracy.

\subsection{Training Procedure}

Training uses the Adam optimizer with \texttt{EarlyStopping} (patience 3, restoring best weights on
validation loss) and \texttt{ReduceLROnPlateau} to anneal the learning rate on validation-loss
plateaus. A custom callback tracks test-set loss at every epoch purely for diagnostic plotting;
critically, this value never influences any training decision, which depends only on the validation
split, preserving the integrity of the held-out test set as a measure of generalization rather than
of implicit tuning.

\section{Standard Classification Results}
\label{sec:standard-results}

Training converged in 19 epochs (batch size 512, 480{,}544 training samples). Table~\ref{tab:biasvar}
summarizes the bias-variance diagnostic.

\begin{table}[h]
\centering
\caption{Bias-variance diagnostic.}
\label{tab:biasvar}
\begin{tabular}{lr}
\toprule
Metric & Value \\
\midrule
5-fold CV validation accuracy & $0.9721 \pm 0.0139$ \\
Train accuracy & 0.9834 \\
Validation accuracy & 0.9827 \\
Test accuracy & 0.9829 \\
Train--test accuracy gap & 0.0005 \\
\bottomrule
\end{tabular}
\end{table}

The negligible train--test gap (0.05 percentage points) combined with a low cross-validation
standard deviation (1.39\%) indicates \textbf{low bias and low variance}: the model is neither
underfitting nor overfitting under the standard evaluation protocol. This diagnostic matters
specifically because it rules out a confound in the LOAO results that follow in
Section~\ref{sec:loao}: the collapse reported there cannot be attributed to the base model being
poorly fit in general, since under the distribution it was actually trained on, it is well-calibrated
by every standard bias-variance measure.

Final test-set performance is summarized in Table~\ref{tab:testresults}.

\begin{table}[h]
\centering
\caption{Test set performance (binary classification, standard split).}
\label{tab:testresults}
\begin{tabular}{lr}
\toprule
Metric & Value \\
\midrule
Accuracy & 0.9829 \\
Precision (weighted) & 0.9829 \\
Recall (weighted) & 0.9829 \\
F1-score (weighted) & 0.9829 \\
Specificity & 0.9782 \\
ROC-AUC & 0.9989 \\
\bottomrule
\end{tabular}
\end{table}

\subsection{Comparative Analysis with State-of-the-Art}
\label{sec:sota}

Table~\ref{tab:sota} places our model against three published results on CICIDS2017.

\begin{table}[h]
\centering
\caption{Comparison against published CICIDS2017 results (standard evaluation only).}
\label{tab:sota}
\small
\begin{tabular}{lrrl}
\toprule
Model & Accuracy & F1 & Parameters \\
\midrule
CNN-LSTM-GRU, parallel fusion \cite{afraji2025} & 0.9949 & 0.9956 & Not reported \\
Hybrid DL, 27 selected features \cite{kumar2023} & 0.9997 & -- & Not reported \\
Benign-only Autoencoder \cite{hindy2020} & 0.75--0.99* & -- & Not reported \\
\textbf{Proposed CNN-LSTM (this work)} & \textbf{0.9829} & \textbf{0.9829} & \textbf{33{,}761} \\
\bottomrule
\end{tabular}
\\[2pt]
\footnotesize{*Autoencoder detection accuracy is category-dependent (75--99\% range across attack
types and thresholds); not a single aggregate figure.}
\end{table}

On standard evaluation, our model does not exceed the highest published figure (99.97\%
\cite{kumar2023}), which uses a curated 27-feature subset and a more elaborate hybrid pipeline. Our
architecture instead prioritizes efficiency: it uses the full feature set with roughly an order of
magnitude fewer parameters than typical published architectures while remaining within the same
accuracy tier (97--99.9\%) reported across CICIDS2017 studies. We state this comparison without
qualification: on this specific metric, under this specific protocol, this work does not claim or
demonstrate superiority over \cite{afraji2025} or \cite{kumar2023}. The contribution of this work
lies instead in Section~\ref{sec:loao}, where none of the compared architectures have published
results at all.

\section{Innovation: Leave-One-Attack-Out Evaluation with Centroid- and Reconstruction-Based Augmentation}
\label{sec:loao}

\subsection{Motivation}

A classifier trained to recognize specific attack categories provides no theoretical guarantee that
it will recognize a category it has never seen, as formalized in Section~\ref{sec:formalization}.
To test this empirically, we implement a Leave-One-Attack-Out (LOAO) protocol: for each of four
attack categories (DoS Hulk, PortScan, DDoS, Bot), we retrain the CNN-LSTM from scratch with that
category entirely excluded from the training pool, then evaluate on a test set containing only the
withheld category plus held-out benign traffic. This directly measures generalization to a
genuinely unseen threat, rather than to a held-out split of an already-seen distribution.

\subsection{Augmentation Method}
\label{sec:augmentation}

Rather than relying solely on the retrained classifier's decision boundary, we test whether two
additional, classifier-independent signals can recover generalization performance.

\paragraph{Centroid-distance features.} The 32-dimensional embedding $g_\theta(x) \in \mathbb{R}^{32}$
from the model's penultimate Dense layer is extracted for every training sample. A benign centroid
and one centroid per \emph{known} attack category (i.e., excluding the withheld one) are computed as
class-wise embedding means:

\begin{equation}
\mu_{\text{benign}} = \frac{1}{|\mathcal{B}|} \sum_{x_i \in \mathcal{B}} g_\theta(x_i), \qquad
\mu_{c} = \frac{1}{|\mathcal{A}_c|} \sum_{x_i \in \mathcal{A}_c} g_\theta(x_i) \;\; \forall c \in \mathcal{C} \setminus \{c_j\}
\end{equation}

where $\mathcal{B}$ is the benign training set and $\mathcal{A}_c$ is the training set for known
attack category $c$. For each query sample, three features are derived: Euclidean distance to the
benign centroid, Euclidean distance to the nearest known-attack centroid, and the signed gap between
the two:

\begin{equation}
d_{\text{benign}}(x) = \lVert g_\theta(x) - \mu_{\text{benign}} \rVert_2, \qquad
d_{\text{attack}}(x) = \min_{c \in \mathcal{C} \setminus \{c_j\}} \lVert g_\theta(x) - \mu_c \rVert_2
\end{equation}

\begin{equation}
\Delta(x) = d_{\text{benign}}(x) - d_{\text{attack}}(x)
\end{equation}

The underlying hypothesis is that a genuinely novel attack should embed far from both known
clusters, or at least measurably closer to the attack region than to benign.

\paragraph{Benign-only autoencoder reconstruction error.} A small autoencoder
(input$\to$32$\to$8$\to$32$\to$input, trained with Adam and mean-squared-error loss) is trained
exclusively on benign traffic in the raw scaled feature space. For encoder $E$ and decoder $D$, the
reconstruction error on a query sample is:

\begin{equation}
r(x) = \frac{1}{m} \sum_{k=1}^{m} \big(x_k - D(E(x))_k\big)^2
\end{equation}

where $m$ is the feature dimensionality. This signal is computed independently of the CNN-LSTM
entirely and is intended as a category-agnostic ``does this resemble normal traffic'' score.

\paragraph{Meta-classifier.} These signals are combined with the base classifier's output
probability $p(x) = f_\theta(x)$ via a lightweight logistic regression meta-classifier, evaluated
under a four-condition ablation so that each component's individual contribution can be isolated:

\begin{align}
\text{Baseline:} \quad & p(x) \ge 0.5 \\
\text{Exp1 (A):} \quad & \sigma(w_0 \, p(x) + b) \\
\text{Exp2 (A+B):} \quad & \sigma\big(w_0 \, p(x) + w_1 \, d_{\text{benign}}(x) + w_2 \, d_{\text{attack}}(x) + w_3 \, \Delta(x) + b\big) \\
\text{Proposed (A+B+C):} \quad & \sigma\big(w_0 \, p(x) + w_1 \, d_{\text{benign}}(x) + w_2 \, d_{\text{attack}}(x) + w_3 \, \Delta(x) + w_4 \, r(x) + b\big)
\end{align}

This design isolates each component's contribution, rather than reporting a single before/after
comparison that would conflate the two signals' individual effects.

\subsection{A Note on Terminology}

This design is inspired by the same intuition behind retrieval-augmented approaches -- using stored
knowledge about known categories to inform a new decision -- but it is \textbf{not}
Retrieval-Augmented Generation (RAG). No generative model is used, and there is no instance-level
retrieval step: distances are computed against fixed, precomputed class centroids rather than
against individually retrieved training examples. A more faithful adaptation of the retrieval
paradigm to this classification setting would replace the centroid comparison with a
$k$-nearest-neighbor search over stored training embeddings, using the label distribution of the
retrieved neighbors as the augmenting signal. We identify this as a concrete direction for future
work (Section~\ref{sec:limitations}) rather than claiming it as already implemented here.

\subsection{Results}
\label{sec:loao-results}

Table~\ref{tab:loao} reports baseline (no augmentation) LOAO performance per withheld category, and
Table~\ref{tab:ablation} reports the full ablation, averaged across all four categories.

\begin{table}[h]
\centering
\caption{Baseline CNN-LSTM performance under LOAO evaluation (category fully withheld from training).}
\label{tab:loao}
\begin{tabular}{lrrr}
\toprule
Withheld Category & Accuracy & Recall & F1 \\
\midrule
DoS Hulk & 0.6319 & 0.6215 & 0.7665 \\
PortScan & 0.7984 & 0.3712 & 0.5088 \\
DDoS & 0.6443 & 0.6313 & 0.7736 \\
Bot & 0.7519 & \textbf{0.0000} & \textbf{0.0000} \\
\bottomrule
\end{tabular}
\end{table}

\begin{table}[h]
\centering
\caption{Ablation study, averaged across all four withheld categories. A = base classifier
probability, B = centroid-distance features, C = autoencoder reconstruction error.}
\label{tab:ablation}
\begin{tabular}{lrrr}
\toprule
Condition & Accuracy & Recall & F1 \\
\midrule
Baseline (A only, direct) & 0.7066 & 0.4060 & 0.5122 \\
Exp1: A (recalibrated) & 0.7063 & 0.4068 & 0.5124 \\
Exp2: A+B & 0.7053 & 0.4063 & 0.5119 \\
Proposed: A+B+C & 0.7053 & 0.4063 & 0.5119 \\
\bottomrule
\end{tabular}
\end{table}

Two findings stand out. First, accuracy collapses substantially relative to the 98.29\%
standard-split result whenever an entire category is withheld, confirming empirically what
Section~\ref{sec:formalization} predicts theoretically: the standard evaluation in
Section~\ref{sec:sota} does not reflect zero-day robustness. Second, and central to the ablation's
purpose, the added centroid-distance and autoencoder signals produce \textbf{negligible change}
relative to the baseline classifier probability alone -- in most conditions the difference is within
0.1--0.3 percentage points, and on Bot specifically, recall remains exactly 0\% across every
configuration tested.

\subsection{Discussion: Interpreting the Negative Result}
\label{sec:discussion}

We do not present the ablation result as an improvement, and we consider it important to state this
plainly rather than overstate the outcome. The value of this experiment lies in what it rules out:
the bottleneck limiting zero-day generalization in this setting is not the \emph{scoring rule}
applied after classification, since neither embedding-distance nor reconstruction-error augmentation
recovers meaningful performance. This points instead to the \emph{representation itself} -- as
argued formally in Section~\ref{sec:formalization}, the 32-dimensional embedding space learned by
the CNN-LSTM is shaped entirely by the attack categories present during training, via the same
optimization (Equation~\ref{eq:erm}) that produces $\theta^{*}$. Distances measured within that
space, even to a well-chosen set of centroids, do not carry reliable information about categories
the space was never shaped to separate, because $g_\theta$ was never optimized with any objective
that would encourage such separation for an unseen class. A reconstruction-error signal computed
independently of the classifier fares no better here, though prior work shows autoencoder-based
detection \emph{can} succeed when attacks are statistically well-separated from benign traffic
\cite{hindy2020} -- our LOAO setting evidently does not provide that separation for the categories
tested.

The complete recall failure on Bot (0.00\% in every condition) is not, in isolation, evidence of a
flaw specific to our implementation. Xu and Liu \cite{xu2025} independently report an identical
0.0000 per-class accuracy for a CNN model on withheld Bot traffic in CICIDS2017, using a different
architecture and a different held-out category set. The convergence of two independent studies on
this exact failure mode is stronger evidence that Bot traffic is intrinsically difficult to
generalize to under a supervised, fully-withheld evaluation than any single result could provide
alone -- and is consistent with the distribution-support argument of Section~\ref{sec:formalization}:
if Bot traffic occupies a region of feature space poorly separated from benign traffic even when
\emph{included} in training (a property independent of which model is used), then no amount of
representational capacity withheld-category retraining can be expected to fix, since the issue is
not a withheld-category artifact but a property of the category's statistical signature relative to
benign traffic. Averaged across all four withheld categories, our baseline LOAO accuracy (70.66\%)
is broadly comparable in magnitude to the unknown-attack accuracy Xu and Liu report for their CNN
model (58.82\%); we note this only as a coarse, order-of-magnitude comparison, since the two studies
withhold different category sets and use different sampling and feature pipelines, not as a
controlled head-to-head benchmark. We do not interpret this comparison as evidence of superiority in
any individual category; where the two studies' withheld categories overlap (Bot), the result is a
tie in total failure, not an advantage for either study.

\section{Real-World Applications}
\label{sec:deployment}

The compact parameter count (33{,}761) makes this architecture a reasonable candidate for
resource-constrained deployment contexts -- edge gateways, IoT hubs, or embedded network appliances
-- where a full-scale ensemble or transformer-based IDS would be impractical. On the
standard-evaluation task, the model is well-suited to complement a signature-based system as a
lightweight secondary classifier for known attack families, operating at low latency with a small
memory footprint.

However, the LOAO results in Section~\ref{sec:loao-results} carry a direct and important deployment
implication: this class of model, on its own, should not be relied upon as a primary defense against
attack categories not represented in its training data. In an enterprise or cloud Security Operations
Center (SOC) context, this motivates a layered deployment strategy -- using the CNN-LSTM classifier
for known-attack triage (where it performs well, Section~\ref{sec:sota}) while retaining human
analyst review or a complementary boundary-based method for traffic flagged as anomalous but not
confidently classified, particularly for categories with a documented history of poor
supervised-model generalization such as Bot and PortScan traffic. The theoretical framing of
Section~\ref{sec:formalization} generalizes beyond this specific architecture: \emph{any} classifier
trained purely by empirical risk minimization over a fixed category set inherits the same
support-limitation argument, suggesting this layered-deployment recommendation applies to the wider
class of supervised NIDS models, not solely to the architecture studied here.

\section{Limitations and Future Work}
\label{sec:limitations}

\begin{itemize}[leftmargin=*]
\item \textbf{Augmentation did not close the generalization gap.} The centroid-distance and
autoencoder signals tested here provided no meaningful recovery of LOAO performance. A concrete
next step, motivated directly by this negative result, is to replace centroid-based comparison with
genuine instance-level $k$-nearest-neighbor retrieval over stored training embeddings, and
separately, to incorporate a contrastive or metric-learning training objective -- for instance, a
triplet or supervised contrastive loss added to Equation~\ref{eq:erm} -- so that the embedding space
itself is shaped to separate ``normal vs.\ anomalous'' rather than ``benign vs.\ these specific known
attacks.''
\item \textbf{Cross-study comparisons are approximate.} The comparison to \cite{xu2025} in
Section~\ref{sec:loao-results} uses a different held-out category set, sampling strategy, and
feature pipeline than this work, and should be read as directional evidence of a shared phenomenon
rather than a controlled benchmark. A fully controlled replication, using an identical withheld
category set and preprocessing pipeline across both architectures, is left to future work.
\item \textbf{Bot detection remains an open problem.} Given its low absolute sample count (1{,}437
records) and the convergent zero-recall finding across two independent studies, few-shot calibration
using a small number of labeled Bot examples -- rather than a fully zero-shot setting -- may be a
more tractable path than further zero-shot augmentation.
\item \textbf{Only four categories were evaluated under LOAO.} A full leave-one-out sweep across
all 14 attack sub-types, while computationally expensive (requiring one full retraining per
category), would give a more complete picture of which categories are and are not tractable under
this protocol, and is a natural extension of this work.
\end{itemize}

\section{Conclusion}
\label{sec:conclusion}

This work builds a parameter-efficient CNN-LSTM intrusion detector that is competitive with
published CICIDS2017 results under standard evaluation, and then subjects it to a stricter, less
commonly reported Leave-One-Attack-Out evaluation that more accurately reflects zero-day deployment
conditions. Consistent with independently published findings \cite{xu2025}, we observe a substantial
and, for one category, total collapse in detection performance under this harder protocol. A
rigorous ablation of two candidate augmentation signals shows neither meaningfully recovers this
lost performance, a negative result that -- grounded in the distribution-shift argument of
Section~\ref{sec:formalization} -- localizes the generalization bottleneck to the learned
representation rather than the post-hoc scoring method, and motivates a specific, actionable
direction -- metric-learning-based representation training and genuine instance-level retrieval --
for future work. We make no claim of outperforming prior work on any individually comparable metric;
the contribution of this work is methodological -- introducing a harder, theoretically motivated
evaluation protocol to a setting where it is largely absent, and providing rigorous, honestly
reported negative evidence about what does and does not help under that protocol.

\begin{thebibliography}{9}

\bibitem{hindy2020}
H. Hindy, R. Atkinson, C. Tachtatzis, J.-N. Colin, E. Bayne, and X. Bellekens,
``Utilising Deep Learning Techniques for Effective Zero-Day Attack Detection,''
\textit{Electronics}, vol. 9, no. 10, p. 1684, 2020.
\url{https://doi.org/10.3390/electronics9101684}

\bibitem{afraji2025}
D. M. A. A. Afraji, J. Lloret, and L. Pe\~nalver,
``An Integrated Hybrid Deep Learning Framework for Intrusion Detection in IoT and IIoT Networks Using
CNN-LSTM-GRU Architecture,''
\textit{Computation}, vol. 13, no. 9, p. 222, 2025.
\url{https://doi.org/10.3390/computation13090222}

\bibitem{kumar2023}
N. Kumar and S. Sharma,
``A Hybrid Modified Deep Learning Architecture for Intrusion Detection System with Optimal Feature
Selection,''
\textit{Electronics}, vol. 12, no. 19, p. 4050, 2023.
\url{https://doi.org/10.3390/electronics12194050}

\bibitem{xu2025}
Z. Xu and Y. Liu,
``Robust Anomaly Detection in Network Traffic: Evaluating Machine Learning Models on CICIDS2017,''
arXiv preprint arXiv:2506.19877, 2025.
\url{https://arxiv.org/abs/2506.19877}

\bibitem{chandola2009}
V. Chandola, A. Banerjee, and V. Kumar,
``Anomaly Detection: A Survey,''
\textit{ACM Computing Surveys}, vol. 41, no. 3, Article 15, 2009.
\url{https://doi.org/10.1145/1541880.1541882}

\bibitem{moustafa2016}
N. Moustafa and J. Slay,
``The Evaluation of Network Anomaly Detection Systems: Statistical Analysis of the UNSW-NB15 Data
Set and the Comparison with the KDD99 Data Set,''
\textit{Information Security Journal: A Global Perspective}, vol. 25, no. 1--3, pp. 18--31, 2016.
\url{https://doi.org/10.1080/19393555.2015.1125974}

\end{thebibliography}

\end{document}
